\chapter{Tiny210
启动与内存初始化}\label{tiny210-ux542fux52a8ux4e0eux5185ux5b58ux521dux59cbux5316}

\section{SD 启动布局}\label{sd-ux542fux52a8ux5e03ux5c40}

Tiny210 不依赖 U-Boot。S5PV210 iROM（BL0）从 SD 卡固定位置加载本仓库的
BL1； BL1 初始化 SDHCI 后读取 MBR，在第一个 FAT32
分区的根目录中查找并加载 \passthrough{\lstinline!xv6\_tiny210.img!}：

{\def\LTcaptype{none} % do not increment counter
\begin{longtable}[]{@{}p{0.18\linewidth}p{0.26\linewidth}p{0.50\linewidth}@{}}
\toprule\noalign{}
SD 扇区 & 内容 & 运行/加载位置 \\
\midrule\noalign{}
\endhead
\bottomrule\noalign{}
\endlastfoot
0 & MBR，制作镜像时保留 & 不执行 \\
1..32 & 16 KiB \passthrough{\lstinline!tiny210-bl1.bin!} & iRAM
\passthrough{\lstinline!0xD0020000!}，代码入口
\passthrough{\lstinline!0xD0020010!} \\
33..48 & 保留 & --- \\
49.. & 原始 uImage 救援副本（可选） & FAT32 文件不可用时才读取 \\
MBR 指定的 FAT32 分区 & 根目录
\passthrough{\lstinline!xv6\_tiny210.img!} & 主启动内核；先放到
DRAM+\passthrough{\lstinline!0x7fc0!}，uImage payload 位于
DRAM+\passthrough{\lstinline!0x8000!} \\
\end{longtable}
}

BL1 的前 16 字节是 iROM 所需的大小、保留字段和字节校验和。链接脚本
\href{../../../kernel/arch/arm/mach-s5pv210/bl1/bl1.ld}{\passthrough{\lstinline!bl1/bl1.ld!}}
强制代码不超过 \passthrough{\lstinline!16 KiB - 16 B!}。

\section{BL0 → BL1 → 内核}\label{bl0--bl1--ux5185ux6838}

\begin{figure}[htbp]
\centering
\resizebox{\linewidth}{!}{%
\begin{tikzpicture}[node distance=7mm]
  \node[hbox] (rom) {iROM / BL0\\读扇区 1--32};
  \node[box,below=of rom] (bl1) {BL1 @ \texttt{0xD0020010}\\校验头、PS\_HOLD、PLL、UART};
  \node[box,below=of bl1] (ddr) {初始化 DMC0 DDR2\\容量和读写测试};
  \node[box,below=of ddr] (sd) {初始化启动通道的 SDHCI PIO\\读取 LBA 0};
  \node[box,below=of sd] (mbr) {校验 MBR \texttt{0x55AA}\\选择首个 \texttt{0x0b/0x0c} 分区};
  \node[box,below left=7mm and 10mm of mbr] (fat) {解析 FAT32 BPB、FAT、根目录\\读取 \texttt{xv6\_tiny210.img}};
  \node[gbox,below right=7mm and 10mm of mbr] (raw) {MBR、FAT32 或文件失败\\扇区 49 救援 uImage};
  \node[box,below=20mm of mbr] (atag) {加载到 \texttt{0x20007fc0} 并校验 uImage\\生成 ATAG：\texttt{r0=0, r1=3466, r2=0x20000100}};
  \node[hbox,below=of atag] (kernel) {跳到 \texttt{0x20008000}\\内核 \texttt{\_entry}};
  \draw[arr] (rom)--(bl1); \draw[arr] (bl1)--(ddr); \draw[arr] (ddr)--(sd); \draw[arr] (sd)--(mbr);
  \draw[arr] (mbr)--(fat); \draw[arr] (mbr)--(raw);
  \draw[arr] (fat)--(atag); \draw[arr] (raw)--(atag); \draw[arr] (atag)--(kernel);
\end{tikzpicture}
}
\caption{MBR/FAT32 内核加载主路径与固定扇区救援回退}
\label{fig:tiny-boot-chain}
\end{figure}


BL1 入口在
\href{../../../kernel/arch/arm/mach-s5pv210/bl1/start.S}{\passthrough{\lstinline!bl1/start.S!}}，
硬件初始化和加载决策在
\href{../../../kernel/arch/arm/mach-s5pv210/bl1/bl1.c}{\passthrough{\lstinline!bl1/bl1.c!}}，
FAT32 只读实现位于
\href{../../../kernel/arch/arm/mach-s5pv210/bl1/bl1fat.c}{\passthrough{\lstinline!bl1/bl1fat.c!}}。
由于实板 iROM 的 \passthrough{\lstinline!CopySDMMCtoMem!} 读取 LBA 0
会卡住，主路径复用
\href{../../../kernel/drivers/mmc/host/sdcore.c}{\passthrough{\lstinline!sdcore.c!}}
的小型 SDHCI PIO 核心读取 MBR、FAT 和文件数据；iROM 复制服务只保留给扇区
49 的救援回退。

BL1 的 FAT32 加载顺序为：

\begin{lstlisting}
LBA 0
  -> 校验 0x55AA
  -> 扫描 4 个 MBR partition entry
  -> 选择第一个 type=0x0b/0x0c 的 FAT32 分区
  -> 读取 VBR/BPB，计算 fat_lba、data_lba、root_cluster
  -> 遍历根目录 cluster chain
  -> 匹配 xv6_tiny210.img 的 VFAT 长文件名或 8.3 别名
  -> 沿 FAT cluster chain 把文件读到 DRAM+0x7fc0
  -> 校验 uImage header、load/entry、文件长度和 data CRC
\end{lstlisting}

只有 MBR/FAT32 无效或根目录找不到内核文件时，BL1 才读取固定扇区 49
的救援 uImage。这样日常升级只需替换 FAT32 文件，不必重新写内核原始扇区。

交接遵循 ARM Linux 传统启动约定：

\begin{lstlisting}
r0 = 0
r1 = 3466 (MINI210 machine ID)
r2 = 0x20000100 (ATAG list)
CPU = SVC mode，IRQ/FIQ 屏蔽，MMU/D-cache 关闭
\end{lstlisting}

\section{\texorpdfstring{\texttt{\_entry}
做什么}{\_entry 做什么}}\label{_entry-ux505aux4ec0ux4e48}

内核入口是
\href{../../../kernel/arch/arm/mach-s5pv210/entry.S}{\passthrough{\lstinline!entry.S!}}
中的 \passthrough{\lstinline!\_entry!}：

\begin{enumerate}
\def\labelenumi{\arabic{enumi}.}
\tightlist
\item
  检查实际运行地址是否等于链接地址。
\item
  拒绝 USR mode 或仍开启 MMU 的错误启动方式。
\item
  关闭 I-cache/D-cache，设置低向量/非 VIC vectored 模式。
\item
  分别设置 IRQ、Abort、Undefined 和 SVC 各模式的栈。
\item
  清空 \passthrough{\lstinline!.bss!}。
\item
  把 VBAR 指向
  \href{../../../kernel/arch/arm/vectors.S}{\passthrough{\lstinline!vectors.S!}}
  的 \passthrough{\lstinline!vectors!}。
\item
  保留 \passthrough{\lstinline!r0/r1/r2!} 并调用
  \passthrough{\lstinline!x210\_main()!}。
\end{enumerate}

链接地址由
\href{../../../kernel/arch/arm/mach-s5pv210/kernel.ld}{\passthrough{\lstinline!kernel.ld!}}
的 \passthrough{\lstinline!KERNEL\_LOAD!} 决定；Tiny210 默认为
\passthrough{\lstinline!0x20008000!}。代码、只读数据、initcall 表、
数据和内嵌 ramdisk 都位于镜像中，\passthrough{\lstinline!.bss!}
不占镜像文件空间。

\section{\texorpdfstring{\texttt{x210\_main()}
内核初始化顺序}{x210\_main() 内核初始化顺序}}\label{x210_main-ux5185ux6838ux521dux59cbux5316ux987aux5e8f}

\href{../../../kernel/arch/arm/mach-s5pv210/setup.c}{\passthrough{\lstinline!setup.c!}}
的 \passthrough{\lstinline!x210\_main()!} 是板级 启动主线：

\begin{lstlisting}
uartinit()                         早期串口可立即打印
read_bootparams(r2)               解析 ATAG 或 DTB
kvminit(); kvminithart()          建立并启用内核页表
kinit1()                          初始化物理页分配器
device_init()                     注册 platform bus
consoleinit(); ttyinit()          控制台与 tty 字符设备
printfinit()
arm32_time_init()                 初始化 tickslock
s5pv210_register_devices()        登记板级控制器资源
do_initcalls()                    驱动注册、匹配并 probe
procinit(); binit(); iinit(); fileinit()
userinit()                        创建第一个用户进程
scheduler()                       不再返回
\end{lstlisting}

文件系统初始化没有直接放在 \passthrough{\lstinline!x210\_main()!}
中，因为日志恢复和 buffer lock 可能睡眠。 第一个普通进程第一次执行
\passthrough{\lstinline!forkret()!} 时调用
\passthrough{\lstinline!fsinit(ROOTDEV)!}，然后才返回用户态。

\section{内存布局与 MMU}\label{ux5185ux5b58ux5e03ux5c40ux4e0e-mmu}

Tiny210 DMC0 的物理 DRAM 从 \passthrough{\lstinline!0x20000000!}
开始。当前端口只把前 128 MiB 交给 页分配器，并在紧随其后的区域保留 2 MiB
framebuffer：

\begin{lstlisting}
0x20000000  +-------------------------------+
            | ATAG @ +0x100                 |
            | kernel @ +0x8000              |
            | text/rodata/data/ramdisk/bss  |
            | kalloc 可用页                  |
0x28000000  +-------------------------------+  DRAM + 128 MiB
            | LCD framebuffer (2 MiB)       |
0x28200000  +-------------------------------+
            | 当前内核未交给页分配器的 DRAM |
0x40000000  +-------------------------------+  512 MiB 顶端
\end{lstlisting}

\href{../../../kernel/arch/arm/vm.c}{\passthrough{\lstinline!vm.c!}}
使用 ARMv7 short-descriptor：

\begin{itemize}
\tightlist
\item
  内核 DRAM 采用 1 MiB section 恒等映射，VA 等于 PA。
\item
  \passthrough{\lstinline!devices.c!} 中每个
  \passthrough{\lstinline!IORESOURCE\_MEM!} 也建立恒等映射。
\item
  控制器寄存器映射为 privileged、strongly ordered、XN。
\item
  framebuffer 映射为 privileged、Normal non-cacheable、XN。
\item
  用户地址范围为 \passthrough{\lstinline![0, 0x20000000)!}，低于 Tiny210
  DRAM 起始地址。
\item
  当前 bring-up 中 CPU cache
  保持关闭，便于先保证设备一致性，但性能较低。
\end{itemize}

\section{第一个用户进程}\label{ux7b2cux4e00ux4e2aux7528ux6237ux8fdbux7a0b}

\passthrough{\lstinline!userinit()!} 的一页用户地址空间里放入
\href{../../../kernel/arch/arm/initcode.S}{\passthrough{\lstinline!initcode.S!}}：

\begin{lstlisting}
userinit()
  -> uvminit(..., initcode)
  -> PC=0, SP=4096, CPSR=user mode
  -> RUNNABLE
  -> scheduler()
  -> forkret()
       -> 首次 fsinit(ROOTDEV)
       -> usertrapret()
  -> initcode: exec("/init")
  -> /init: 创建 /dev 节点、连接 ttyS0、fork/exec /bin/sh
\end{lstlisting}

Tiny210 构建实际使用
\href{../../../kernel/arch/arm/init.c}{\passthrough{\lstinline!kernel/arch/arm/init.c!}}，
而不是仓库根用户集合中的 \passthrough{\lstinline!user/init.c!}。

\section{启动故障定位}\label{ux542fux52a8ux6545ux969cux5b9aux4f4d}

{\def\LTcaptype{none} % do not increment counter
\begin{longtable}[]{@{}p{0.34\linewidth}p{0.60\linewidth}@{}}
\toprule\noalign{}
最后可见信息 & 优先检查 \\
\midrule\noalign{}
\endhead
\bottomrule\noalign{}
\endlastfoot
完全无串口 & SD 启动拨码、BL1 扇区、UART0/COM0、波特率 \\
\passthrough{\lstinline!BL1: ... halted!} &
PLL、DDR、SDHCI、MBR/FAT32、uImage CRC \\
\passthrough{\lstinline!\_entry!} early error &
加载地址、\passthrough{\lstinline!bootm!} 交接、MMU 状态 \\
打印内核 banner 后停止 & MMU 映射、kalloc、platform device/initcall \\
timer 日志前停止 & VIC 或 PWM 资源/时钟 \\
\passthrough{\lstinline!panic: invalid file system!} & 内嵌
\passthrough{\lstinline!fs.img!}、superblock 或构建目录混用 \\
\passthrough{\lstinline!/init!} 或 shell 找不到 &
\passthrough{\lstinline!mkfs!} 输入程序和 \passthrough{\lstinline!/bin!}
目录内容 \\
\end{longtable}
}
